% Generated draft: requires mathematical and source review.
\documentclass[11pt]{article}
\usepackage[margin=25mm]{geometry}
\usepackage{fontspec}
\setmainfont{Latin Modern Roman}
\usepackage{amsmath,amssymb,mathtools,unicode-math}
\setmathfont{Latin Modern Math}
\usepackage{enumitem,needspace,xurl,hyperref}
\setlength{\parindent}{0pt}
\setlength{\parskip}{5pt}
\setlength{\emergencystretch}{3em}
\newcommand{\N}{\mathbb{N}}
\newcommand{\Z}{\mathbb{Z}}
\newcommand{\Q}{\mathbb{Q}}
\newcommand{\R}{\mathbb{R}}
\newcommand{\C}{\mathbb{C}}
\newcommand{\F}{\mathbb{F}}
\newcommand{\A}{\mathbb{A}}
\newcommand{\PP}{\mathbb{P}}
\newcommand{\E}{\mathbb{E}}
\newcommand{\eps}{\varepsilon}
\newcommand{\dd}{\,\mathrm{d}}
\begin{document}
\section*{UnsolvedMath ID: 700194}
\textbf{Source problem code:} catalogue entry
\subsection{Definitions and mathematical statement}
A group \(G\) is \emph{torsion-free} if, for every \(g\in G\) and every integer \(n\geq 1\), the equality \(g^n=1\) implies \(g=1\). It is \emph{locally nilpotent} if every finitely generated subgroup is nilpotent. Here a group \(H\) is nilpotent if its lower central series, defined by \(\gamma_1(H)=H\) and \(\gamma_{k+1}(H)=[\gamma_k(H),H]\), reaches the trivial subgroup after finitely many steps; \([A,B]\) denotes the subgroup generated by the commutators \(a^{-1}b^{-1}ab\) with \(a\in A\) and \(b\in B\).

A subgroup \(K\leq G\) is \emph{isolated in \(G\)} if
\[
\forall g\in G\ \forall n\in\Z_{\geq 1},\qquad g^n\in K\ \Longrightarrow\ g\in K.
\]
Let \(\mathfrak{X}\) be a hereditary radical of groups, and let \(\mathfrak{X}(G)\) denote its associated radical subgroup of \(G\). The supplied record does not specify the axioms for a hereditary radical or the construction of \(\mathfrak{X}(G)\). These terms are therefore retained in the sense intended by the cited source; the quantification over radicals below requires that convention to be fixed before this is a fully self-contained formulation.

\textbf{Question (B. I. Plotkin, 1969).} For every hereditary radical \(\mathfrak{X}\) and every locally nilpotent torsion-free group \(G\), is \(\mathfrak{X}(G)\) isolated in \(G\)? Equivalently, under the intended radical convention, must
\[
\forall g\in G\ \forall n\in\Z_{\geq 1},\qquad
 g^n\in\mathfrak{X}(G)\ \Longrightarrow\ g\in\mathfrak{X}(G)
\]
hold?
\subsection{Short English statement}
Plotkin asks whether every hereditary radical subgroup of a locally nilpotent torsion-free group contains every element of the ambient group whose positive power it contains. The source labels the problem open; that label is supplied metadata, not an independently verified assessment. The intended definition of hereditary radical requires source clarification.
\subsection{Sources}
\begin{itemize}
\item Supplied local UnsolvedMath registry record: permanent ID \(700194\), title \emph{Kourovka Notebook Problem 3.45}, exact source problem code \texttt{catalogue entry}; upstream revision \texttt{d371ee917ac9647fbb06cd6fc529deec6f3dda82}. The record was created and updated on \texttt{2026-10-16T14:00:00Z} and supplies the status \texttt{open}.
\item The record attributes the problem to B. I. Plotkin, proposed in 1969, and cites \emph{The Kourovka Notebook (Unsolved Problems in Group Theory)}, 3rd issue (1969), problem 3.45, as included in the 21st-issue edition (2026), \texttt{arXiv:1401.0300v48}.
\item The supplied collection description identifies E. I. Khukhro and V. D. Mazurov as editors of the Kourovka Notebook collection covering Issues 1--21 (1965--2026). The bibliographic attribution and edition information above are supplied source data and have not been independently verified.
\end{itemize}
\end{document}
